\section{Source realization of the enriched Stein--Lyapunov state}
\label{sec:ns-owner-closure}

The construction in the \cref{sec:ns-pcf-prospective-construction} brings the complete PC--Fock tower, bidirectional memory, carry, microscopic innovation, spectral channel, homogeneous force coordinate, unity, and exceptional publication into a single state space. Physical time parametrizes terminal storage, while nonadic depth parametrizes the Stein recurrence. The construction governing this section is the \texttt{OWNER\_GOVERNING} of the enriched HMT state. KYP, NS--OBS, and LRDN--LCN charts are retained only as \texttt{CONTROL\_NON\_GOVERNING} of reduced representations.

\subsection{PC--Fock colligation and telescoping}

Let
\[
 \mathbf x_{M,j}^{\rm own}
 \in\mathscr K_{M,j}^{\rm PCF}
 \widehat\otimes\mathcal E_{\rm exc}
 \oplus\mathscr G_M^{F,[j]}
\]
the homogeneous state of the physical diagonal. For every \(M,J,T\), the source metric is constructed prospectively as the Gram matrix of the terminal column and all future losses of that same diagonal:
\begin{equation}
 \begin{aligned}
 \left\langle x,U_{M,j}^{\rm own}x\right\rangle
 :={}&
 \mathbb E_{j:J}
 \left\|\mathcal E_{M,J,T}^{\rm term}
       \Gamma_{M,j:J}^{\rm own}x\right\|^2\\
 &+\sum_{k=j}^{J-1}\mathbb E_{j:k}
 \left\|L_{M,k,T}^{\rm own}
       \Gamma_{M,j:k}^{\rm own}x\right\|^2 .
 \end{aligned}
 \label{eq:nsclose-owner-metric-definition}
\end{equation}
Every term is a square of coordinates already constructed: viscous degree one, convective degree two, carry, innovation, spectral archive, memory \(D/H\), and homogeneous force. Separating the first step of the sum gives, before introducing any KYP root,
\begin{equation}
 \boxed{
 U_{M,j}^{\rm own}
 =
 (\Gamma_{M,j}^{\rm own})^*
 U_{M,j+1}^{\rm own}\Gamma_{M,j}^{\rm own}
 +(L_{M,j,T}^{\rm own})^*L_{M,j,T}^{\rm own},
 \qquad
 G_{M,J,T}^{\rm own}\preceq C_TU_{M,J}^{\rm own}.}
 \label{eq:nsclose-proprietary-terminal-gate}
\end{equation}
\label{eq:nsclose-owner-identity}
In the terminal normalization of \eqref{eq:ns-import-terminal-column}, one may take \(C_T=1\). The equality of maps, amplified by the identity in the Moonshine layer, produces the isometric cascade
\begin{equation}
 (\mathfrak U_{M,J,T}^{\rm own})^{\sharp}
 \mathfrak U_{M,J,T}^{\rm own}=I.
 \label{eq:nsclose-proprietary-full-cascade}
\end{equation}
Iterating \eqref{eq:nsclose-owner-identity} yields the telescoping identity
\begin{equation}
 U_{M,0}^{\rm own}
 =
 (\Gamma_{M,0:J}^{\rm own})^*U_{M,J}^{\rm own}
  \Gamma_{M,0:J}^{\rm own}
 +\sum_{j<J}(\Gamma_{M,0:j}^{\rm own})^*
  (L_{M,j,T}^{\rm own})^*L_{M,j,T}^{\rm own}
  \Gamma_{M,0:j}^{\rm own}.
 \label{eq:nsclose-owner-telescope}
\end{equation}
Its orthogonal balance ledger is
\begin{equation}
 \boxed{
 E_s(u_M(t))\le\|Z_{M,J,t}\|^2
 =\mathbb E_{0:J}G_{M,J,T}(t)
 +\sum_{j<J}\mathbb E_{0:j}
   \Lambda_{M,j}^{\rm own}([0,t]).}
 \label{eq:nsclose-proprietary-ledger}
\end{equation}
The common root \(\Xi_{M,0}=J_{0\to M}\Xi_0\), evolutionary preservation of unity, and the carry port of \eqref{eq:ns-import-carry-port} ensure that the right-hand side is charged only once. The source budget is
\begin{equation}
 \begin{aligned}
 \mathcal B_T^{\rm own}:={}&
 C_{\rm HMT}
 +(1+2\widetilde M_{s,\beta,\nu}^2)
  e^{R_{\rm F}^2\|u_0\|_{H^s}^2}\\
 &+C_{\nu,s}\int_0^T\|f(t)\|_{H^{s-1}}^2\,dt
 +\|\mathcal J_T^{\rm own}(u_0,f)\|^2 ,
 \end{aligned}
 \label{eq:nsclose-owner-data-budget}
\end{equation}
where \(\mathcal J_T^{\rm own}\) is the root injection determined by the initial state transported by TPK, its two sheets, and the forced input. Its amplifications are
\[
 \mathcal J_{M,J,T}^{\rm own}
 :=I_{0\to M,J}\mathcal J_T^{\rm own},
 \qquad I_{0\to M,J}^*I_{0\to M,J}=I.
\]
By \eqref{eq:nsclose-owner-telescope}, this is not a quantity reconstructed from the future orbit. Preservation of unity and coisometry of the two-way reader give, on the declared domain,
\begin{equation}
 \sup_{M,J}\|\mathcal J_{M,J,T}^{\rm own}(u_0,f)\|^2
 \le C_{{\rm HMT},T}
 \left(1+\|u_0\|_{H^{s+1}}^2
 +\int_0^T\|f(t)\|_{H^{s-1}}^2\,dt\right).
 \label{eq:nsclose-owner-root-uniformity}
\end{equation}
Thus \(\mathcal B_T^{\rm own}\) is independent of \(M,J\). The ledger gives
\begin{equation}
 \boxed{
 \sup_{M,J}\left[
  \sup_{0\le t\le T}\|Z_{M,J,t}\|^2
  +\frac\nu2\int_0^TD_{s,M}(t)\,dt
 \right]\le\mathcal B_T^{\rm own}<\infty.}
 \label{eq:nsclose-proprietary-uniform-budget}
\end{equation}

The hydrodynamic publication is the degree-one return of the same cascade:
\begin{equation}
 R_{M,J,T}^{\rm full}
 :=\mathcal R_{1,M}
   (\mathfrak U_{M,J,T}^{\rm own})^{\sharp},
 \qquad
 \boxed{
 R_{M,J,T}^{\rm full}Z_{M,J,t}=u_M(t).}
 \label{eq:nsclose-proprietary-exact-return}
\end{equation}
\label{eq:nsclose-exact-full-return}
This equality coincides with \eqref{eq:ns-import-exact-return} and avoids evaluating an individual refinement branch.

\begin{theorem}[Uniform estimate for the source PC--Fock realization]
\label{thm:nsclose-proprietary-full}
For every \(T<\infty\), the Galerkin approximations satisfy
\begin{equation}
 \sup_M\sup_{0\le t\le T}\|u_M(t)\|_{H^s}^2
 +\nu\sup_M\int_0^T\|u_M(t)\|_{H^{s+1}}^2\,dt
 \le C_s\mathcal B_T^{\rm own},
 \label{eq:nsclose-uniform-sobolev-precompactness}
\end{equation}
\label{eq:nsclose-uniform-sobolev}
with a constant independent of \(M\) and of nonadic depth.
\end{theorem}

\begin{proof}
Identity \eqref{eq:nsclose-proprietary-exact-return} and contractivity of the return give \(\|u_M(t)\|_{H^s}^2\le C_s\|Z_{M,J,t}\|^2\). Equivalence of \(D_{s,M}\) with the \(H^{s+1}\) norm on the mean-zero solenoidal subspace, together with \eqref{eq:nsclose-proprietary-uniform-budget}, proves the estimate.
\end{proof}

\subsection{Alternative control through a route metric}

The following chart checks refinement in normalized coordinates. It neither constructs nor conditions \(U^{\rm own}\): it omits unity, exceptional publication, and part of the affine channel already present in \eqref{eq:nsclose-owner-metric-definition}.

The Ambivalence chart supplies
\begin{equation}
 G_{\rm av}=\begin{pmatrix}2&-1\\-1&1\end{pmatrix},
 \qquad
 M_{\rm av}=\begin{pmatrix}3&-1\\1&0\end{pmatrix},
 \qquad
 M_{\rm av}^*G_{\rm av}M_{\rm av}
 \preceq\varphi_{\rm HMT}^4G_{\rm av}.
 \label{eq:nsclose-ambivalence-route-metric}
\end{equation}
After normalization \(\widehat M_j=M_j/3\),
\begin{equation}
 \mathsf R_j:=G_j-\widehat M_j^*G_{j+1}\widehat M_j\succeq0.
 \label{eq:nsclose-route-exact-defect}
\end{equation}
In the stationary chart,
\[
 \mathsf R_{\rm av}
 =\frac19\begin{pmatrix}5&-4\\-4&7\end{pmatrix},
 \qquad
 \lambda_{\min}(\mathsf R_{\rm av})
 =\frac{6-\sqrt{17}}9.
\]
With the graph metric
\[
 \mathsf G_M^{\rm gr}
 =I+(\mathcal Y_{\beta,M})^*\mathcal Y_{\beta,M},
 \qquad
 \mathcal C_M
 =\mathcal Y_{\beta,M}(\mathsf G_M^{\rm gr})^{-1/2},
\]
we define
\begin{equation}
 \mathcal A_{M,j}=\widehat M_j\otimes I,
 \qquad
 \mathcal O_{M,j}=\mathsf R_j^{1/2}\otimes\mathcal C_M,
 \qquad
 \mathcal D_{M,j}
 =\mathsf R_j^{1/2}\otimes(I-\mathcal C_M^*\mathcal C_M)^{1/2}.
 \label{eq:nsclose-owner-route-column}
\end{equation}
The route Stein identity is
\begin{equation}
 \boxed{
 \mathsf U_{M,j}^{\rm PCF}
 -\mathcal A_{M,j}^*\mathsf U_{M,j+1}^{\rm PCF}\mathcal A_{M,j}
 =\mathcal O_{M,j}^*\mathcal O_{M,j}
  +\mathcal D_{M,j}^*\mathcal D_{M,j}.}
 \label{eq:nsclose-owner-stein}
\end{equation}
Its iteration gives
\begin{equation}
 \begin{aligned}
 \mathsf U_{M,0}^{\rm PCF}
 -\mathcal A_{M,0:J}^*\mathsf U_{M,J}^{\rm PCF}\mathcal A_{M,0:J}
 =\sum_{j<J}\mathcal A_{M,0:j}^*
 \bigl(\mathcal O_{M,j}^*\mathcal O_{M,j}
      +\mathcal D_{M,j}^*\mathcal D_{M,j}\bigr)
 \mathcal A_{M,0:j}.
 \end{aligned}
 \label{eq:nsclose-owner-stein-telescope}
\end{equation}
This factorization describes refinement geometry. Affine recurrence \eqref{eq:ns-import-affine-stein} also incorporates physical evolution, Carleman cross terms, and the force.

Bidirectional memory turns route loss into a budget. With \(\mathcal N_{M,j}=2(r-1)M_{M,j}^-\),
\begin{equation}
 E_j\Delta_T\mathcal N_{M,j+1}
 =\Delta_T\mathcal N_{M,j}+2K_{M,j}(T),
 \qquad
 K_{M,j}(T)=\int_0^T\kappa_{M,j}(t)\,dt.
 \label{eq:nsclose-carry-memory-telescope}
\end{equation}
Telescoping and the root bound yield
\begin{equation}
 \mathbb E_{0:J}G_{M,J,T}(0)
 \le E_{s,M}(P_Mu_0)+|\Delta_T\mathcal N_0|
 +\|\Xi_{M,0}\|^2+Q_f(T),
 \label{eq:nsclose-owner-terminal-bound}
\end{equation}
which is the abbreviated form of route control \eqref{eq:ns-import-root-bound}--\eqref{eq:ns-import-explicit-budget}. This chart records the cost of the negative sheet,
\[
 \Delta_T\mathcal N_{M,0}
 =2\int_0^T B^-_{s,M,0}(u_M(t))\,dt.
\]
In this reduced representation, the inequality
\[
 X^-_{M,T}(0)\preceq C_T\mathsf U^{\rm PCF}_{M,0},
 \qquad \sup_M C_T<\infty
 \tag{NS--OBS}
\]
is a useful falsifier of the chart. Its failure does not transfer to the source construction, because \(\Delta_T\mathcal N_{M,0}\) does not replace root injection \(\mathcal J_T^{\rm own}\) of \eqref{eq:nsclose-owner-data-budget}--\eqref{eq:nsclose-owner-root-uniformity}. Similarly, KYP and LRDN--LCN verify completions or dominations of reduced readers; none is a premise of the \cref{thm:nsclose-proprietary-full}.

\subsection{Existence, uniqueness, and global smoothness}

\begin{theorem}[Global regularity of Navier--Stokes]
\label{thm:nsclose-global-smooth}
\label{thm:realizaciones-ns}
Let \(s>5/2\), \(\nu>0\), let \(u_0\in H^{s+1}_\sigma(\mathbb T^3)\) have mean zero, and let \(f\in L^2_{\rm loc}([0,\infty);H^{s-1}_\sigma)\). There exists a unique global solution
\begin{equation}
 u\in L^\infty_{\rm loc}([0,\infty);H^s_\sigma)
 \cap L^2_{\rm loc}([0,\infty);H^{s+1}_\sigma),
 \qquad
 \partial_tu\in L^2_{\rm loc}([0,\infty);H^{s-1}_\sigma).
 \label{eq:nsclose-global-regularity-class}
\end{equation}
For smooth data and force, \(u\) and the normalized pressure are smooth.
\end{theorem}

\begin{proof}
The \cref{thm:nsclose-proprietary-full} gives a uniform bound for \((u_M)_M\) in \(L^\infty(0,T;H^s)\cap L^2(0,T;H^{s+1})\). The Galerkin equation and continuity
\[
 H^s\times H^s\longrightarrow H^{s-1},
 \qquad (u,v)\longmapsto\mathbb P(u\cdot\nabla v),
\]
bound \(\partial_tu_M\) in \(L^2(0,T;H^{s-1})\). Since \(H^{s+1}\Subset H^s\hookrightarrow H^{s-1}\), Aubin--Lions yields, after extracting a subsequence,
\[
 u_M\longrightarrow u\quad\hbox{strongly in }L^2(0,T;H^s).
\]
This convergence permits passage to the limit in \(\mathbb P(u_M\cdot\nabla u_M)\), while weak convergence preserves the viscous term and initial datum. A strong solution in \([0,T]\) is obtained.

If \(u,v\) are two solutions and \(w=u-v\), the product estimate and Young's inequality give
\[
 \frac12\frac d{dt}\|w\|_{H^{s-1}}^2
 +\frac\nu2\|w\|_{H^s}^2
 \le C_{s,\nu}\bigl(\|u\|_{H^s}^2+\|v\|_{H^s}^2\bigr)
 \|w\|_{H^{s-1}}^2.
\]
Grönwall proves uniqueness. Repeating the construction on horizons \(T=1,2,\ldots\), the solutions agree on overlapping intervals and yield the global solution.

The pressure is recovered from
\begin{equation}
 -\Delta p=\partial_i\partial_j(u_i u_j)-\operatorname{div}f,
 \qquad \int_{\mathbb T^3}p\,dx=0.
 \label{eq:nsclose-pressure-recovery}
\end{equation}
Iterating the estimate in \(s+k\) and the equation recovers all spatial and temporal derivatives when \(u_0\) and \(f\) are smooth.
\end{proof}

The structural contribution consists in representing convection as a linear output of the PC--Fock tower, preserving it in memory \(D/H\), and incorporating it into an orthogonal identity with exact physical return. The resulting uniform estimate is then published as global regularity of the hydrodynamic equation.
